% TOP_PROBLEM: {"id": 3397, "title": "The robustness of the tensor product", "category_id": 15, "category": {"id": 15, "name": "computer_science", "display_name": "Computer Science", "description": "Computational complexity, algorithms, and theoretical CS.", "slug": "computer-science", "order_index": 15, "created_at": "2026-07-31T15:26:25.671Z"}, "status": "open", "problem_number": "OPG-445", "set_id": 12}
% FIRST_POSTED: 2026-10-02
% ATTEMPT_STATUS: unresolved
% UnsolvedMath ID 3397; source catalogue code OPG-445
% !TeX program = lualatex
% Research attempt by GPT-6 Astra Ultra; no claim of a new complete solution.
\documentclass[11pt,a4paper]{article}
\usepackage[margin=25mm]{geometry}
\usepackage{fontspec}
\setmainfont{lmroman10-regular.otf}[BoldFont=lmroman10-bold.otf,
 ItalicFont=lmroman10-italic.otf,BoldItalicFont=lmroman10-bolditalic.otf]
\usepackage{amsmath,amssymb,mathtools}
\usepackage{unicode-math}
\setmathfont{latinmodern-math.otf}
\AtBeginDocument{\let\setminus\smallsetminus}
\usepackage{xurl,hyperref}
\hypersetup{colorlinks=true,urlcolor=blue}
\setcounter{secnumdepth}{0}
\setlength{\parindent}{0pt}
\setlength{\parskip}{5pt}
\setlength{\emergencystretch}{3em}
\begin{document}
\section*{The robustness of the tensor product}\label{3397}
{\small UnsolvedMath ID 3397 · OPG-445 · Computer Science · OpenGarden}
\subsection{Definitions and mathematical statement}
Fix a finite alphabet $A$. Let $R\subseteq A^n$ and $C\subseteq A^m$
be nonempty codes whose product
$P=R\otimes C\subseteq A^{m\times n}$ is nonempty: every row of
a matrix in $P$ belongs to $R$, and every column belongs to $C$.
For words of equal length, relative Hamming distance is the proportion
of differing coordinates; distance to a code is its minimum over
codewords. Use the analogous normalized entry distance for matrices.
For $M\in A^{m\times n}$, put
\[
 \rho(M)=\frac12\left(\frac1m\sum_{i=1}^m d(M_{i,*},R)
                  +\frac1n\sum_{j=1}^n d(M_{*,j},C)\right).
\]
The pair is $\alpha$-robust if $\rho(M)\ge\alpha d(M,P)$ for every
$M$. Characterize families $(R_t,C_t)$ for which one constant
$\alpha>0$ works for all $t$ as the lengths grow.

The family quantifier is essential: every fixed nonempty finite
product has some positive minimum ratio over matrices outside it.
The desired bound is independent of the lengths. The source also
allows the especially important case of linear codes over a finite
field; the definition above does not require linearity. Characterize
the component-code conditions that yield a uniform lower bound,
using this average row-and-column test and this precise distance.
\subsection{Short English statement}
Which families of code pairs make average row and column errors control the distance of a whole matrix from their product code by one fixed constant?
\subsection{Research attempt}
\paragraph{Scope and process.}
This is a GPT-6 Astra Ultra research attempt dated 2 October 2026, using
primary-source browsing and elementary mathematical checks. The canonical
definition above is retained. Results below are restricted results or
reductions, not a claim of a new solution to the entire catalogue question.
No formal proof assistant or external mathematical referee checked this text.


\paragraph{Literature and scope.}
Goldreich and Meir [R1] construct good component codes whose tensor
product is not robustly testable: positive rate and relative distance
alone are insufficient. This is a warning against claiming a general
characterization from those two parameters. The canonical question
also allows nonlinear codes and does not require positive rate.
We prove two elementary uniform robustness statements in that full
setting and identify why a finite verification cannot settle a family.
All distances below are normalized by the number of entries concerned.

\paragraph{Uniform robustness for repetition codes.}
Take $R=\{a^n:a\in A\}$ and $C=\{a^m:a\in A\}$. Their product consists
exactly of the constant matrices. Given $M$, replace each row by a
nearest constant row and call the resulting matrix $U$. Write
$r=d(M,U)$, the mean row distance. Independently replace each column
by a nearest constant column to form $V$, and write $c=d(M,V)$.
Ties can be resolved arbitrarily. The triangle inequality gives
\[
                   \delta=d(U,V)\le r+c.
\]
Let $a_i$ be the constant symbol in row $i$ of $U$, and $b_j$ the
symbol in column $j$ of $V$. If $p_a$ and $q_a$ are the proportions
of row and column labels equal to $a$, then
\[
                   1-\delta=\sum_{a\in A}p_aq_a
                              \le\max_{a\in A}p_a.
\]
Choose a symbol attaining that maximum. The distance from $U$ to the
all-$a$ matrix is $1-p_a\le\delta$, so
$d(M,P)\le r+\delta\le2r+c$. Interchanging rows and columns gives
$d(M,P)\le r+2c$. Taking the smaller bound yields
\[
 d(M,P)\le\min(2r+c,r+2c)
          \le\tfrac32(r+c)=3\rho(M).
\]
Thus $\alpha=1/3$ works uniformly for all positive $m,n$ and every
finite alphabet. No linearity or field operations were used.

This constant is a certificate, not a claim of optimality. Over the
binary alphabet, take an even number of rows, half constantly zero
and half constantly one. Then $r=0$, $c=1/2$ and $d(M,P)=1/2$.
Any universal constant for repetition-code pairs is consequently at
most $1/2$. The elementary argument confines that best constant to
$[1/3,1/2]$ without deciding its exact value.

\paragraph{A second family and a failed necessary condition.}
If $R=A^n$ is the full row code and $C$ is any nonempty code, columns
can be repaired independently to their nearest codewords. There are
no further row restrictions. Therefore
\[
 d(M,R\otimes C)=\frac1n\sum_j d(M_{*,j},C),
       \qquad \rho(M)=\tfrac12d(M,R\otimes C).
\]
Every such family is uniformly $1/2$-robust. For a proper product this
constant is exact; when the product is the whole matrix space the
inequality is vacuous. In particular, over $A=\{0,1\}$ let
$C_m=\{0^m,10^{m-1}\}$ and let $R_n=A^n$. Both are linear binary codes,
but their minimum relative distances are $1/m$ and $1/n$, tending
to zero. The product remains $1/2$-robust as both dimensions grow.
Thus a uniform positive lower bound on the distances of both components
is not necessary for the broad characterization stated here. This
example does not contradict results that impose additional rate or
nontriviality hypotheses.

\paragraph{Verification and the family quantifier.}
Certificate [R2] enumerates every binary matrix for all dimensions
$1\le m,n\le3$, computes the row, column and constant-matrix distances
exactly, and checks the $1/3$ inequality. It also checks the striped
example above. These finite computations test normalizations and edge
cases; the proof supplies the constant independently of dimensions.
Indeed, for any fixed nonempty finite product, a matrix outside it
has at least one invalid row or column. Correcting that line needs
at least one entry change, giving $\rho(M)\ge1/(2mn)$. Since
$d(M,P)\le1$, the bound $\rho(M)\ge d(M,P)/(2mn)$ follows.
This establishes positivity for each fixed pair but tends to zero
with the lengths and says nothing about a uniform family constant.

\paragraph{Final status and first remaining gap.}
\textbf{UNRESOLVED.} Two broad special families admit uniform
constants, and positive component distance is shown unnecessary in
the stated generality. A characterization for arbitrary code pairs
still requires control of the global compatibility of row and column
repairs. The repetition proof obtains that control from constant labels;
there is no replacement argument here for unrestricted codes.

\subsection{Sources}
\begin{itemize}
\item[S1] ULAM.AI and the UnsolvedMath Contributors, supplied problems.json, record 3397 (OPG-445), OpenGarden collection. Catalogue identity and source transcription; CC BY 4.0.
\url{https://huggingface.co/datasets/ulamai/UnsolvedMath}.
\item[S2] Open Problem Garden, The robustness of the tensor product. Original problem and discussion; the mathematical formulation below is expanded from the supplied source record.
\url{http://www.openproblemgarden.org/op/the_robustness_of_the_tensor_product_0}.

\item[R1] Oded Goldreich and Or Meir, \emph{The tensor product of two
good codes is not necessarily robustly testable}, Information Processing
Letters 112 (2012), 351--355. Primary ECCC revision, including Theorem 3.1, checked
2 October 2026.
\url{https://eccc.weizmann.ac.il/report/2007/062/revision/2/download}.
\item[R2] Reproducible exact checks:
\path{scripts/check_attempt_batch03_colouring_2026_10_02.py}, section 3397.

\end{itemize}
\end{document}
